\chapter{Rechenregeln für Schrankenmengen}
\label{b13guide:ch:bounds}

\section{Elemente, Teilmengen und Randfälle}

\begin{theorem}[Schranken als Einschließungen]
Für $S\subseteq P$ und $\ell,u\in P$ gilt
\[
 \ell\in\Lo(S)\ \Longleftrightarrow\ S\subseteq\mathord\uparrow\ell,
 \qquad
 u\in\Hi(S)\ \Longleftrightarrow\ S\subseteq\mathord\downarrow u.
\]
Für $p\in P$ folgen aus $p\leq\ell\in\Lo(S)$ die Zugehörigkeit
$p\in\Lo(S)$ und aus $u\in\Hi(S)$ sowie $u\leq p$ die
Zugehörigkeit $p\in\Hi(S)$.
\end{theorem}

\begin{theorem}[Umkehrung von Inklusionen]\label{b13guide:thm:antitone}
Für $S\subseteq T\subseteq P$ gilt
\[
 \Lo(T)\subseteq\Lo(S),\qquad \Hi(T)\subseteq\Hi(S).
\]
Jede nach oben beziehungsweise unten beschränkte Menge vererbt
diese Eigenschaft auf alle ihre Teilmengen.
\end{theorem}

\begin{theorem}[Leere und einelementige Mengen]
\label{b13guide:thm:empty-bounds}
Es gilt
\[
 \Lo(\varnothing)=\Hi(\varnothing)=P.
\]
Die leere Menge ist genau dann nach oben beschränkt, wenn $P\ne\varnothing$;
dieselbe Bedingung gilt für die Beschränktheit nach unten.
Für jedes $p\in P$ gilt
\[
 \Lo(\{p\})=\mathord\downarrow p,\qquad
 \Hi(\{p\})=\mathord\uparrow p.
\]
Insbesondere ist jede einelementige Teilmenge beschränkt.
\end{theorem}

\begin{theorem}[Vergleich unterer und oberer Schranken]
\label{b13guide:thm:cross-bounds}
Ist $S\ne\varnothing$, so gilt
\[
 \ell\in\Lo(S),\ u\in\Hi(S)\quad\Rightarrow\quad\ell\leq u.
\]
Allgemeiner gilt für $S\cap T\ne\varnothing$
\[
 \ell\in\Lo(S),\ u\in\Hi(T)\quad\Rightarrow\quad\ell\leq u.
\]
Ohne die jeweilige Nichtleerheitsvoraussetzung gilt die Folgerung
im allgemeinen nicht.
\end{theorem}


\begin{theorem}[Beschränktheit und einschließende Intervalle]
\label{b13guide:thm:bounded-interval}
Schreibe $[a,b]_P=\{x\in P\mid a\leq x\leq b\}$ für $a,b\in P$.
Für jede Teilmenge $S\subseteq P$ gilt
\[
\begin{aligned}
 S\text{ ist nach oben beschränkt}
   &\ \Longleftrightarrow\ \exists b\in P\;(S\subseteq\mathord\downarrow b),\\
 S\text{ ist nach unten beschränkt}
   &\ \Longleftrightarrow\ \exists a\in P\;(S\subseteq\mathord\uparrow a),\\
 S\text{ ist beschränkt}
   &\ \Longleftrightarrow\ \exists a,b\in P\;(S\subseteq[a,b]_P).
\end{aligned}
\]
Für $S\ne\varnothing$ erfüllen solche Endpunkte stets $a\leq b$.
Die Äquivalenzen gelten auch für $S=\varnothing$;
bei $P=\varnothing$ sind ihre beiden Seiten jeweils falsch.
\end{theorem}

\section{Mengenoperationen}

\begin{theorem}[Vereinigung, Durchschnitt und Hinzufügen eines Elements]
\label{b13guide:thm:operations}
Für $S,T\subseteq P$ und $p\in P$ gilt
\[
\begin{aligned}
 \Lo(S\cup T)&=\Lo(S)\cap\Lo(T),\\
 \Hi(S\cup T)&=\Hi(S)\cap\Hi(T),\\
 \Lo(S)\cup\Lo(T)&\subseteq\Lo(S\cap T),\\
 \Hi(S)\cup\Hi(T)&\subseteq\Hi(S\cap T),\\
 \Lo(S\cup\{p\})&=\Lo(S)\cap\mathord\downarrow p,\\
 \Hi(S\cup\{p\})&=\Hi(S)\cap\mathord\uparrow p.
\end{aligned}
\]
Bei den Durchschnittsformeln kann die Inklusion echt sein.
Die Vereinigung ist genau dann nach oben beschränkt, wenn
$\Hi(S)\cap\Hi(T)\ne\varnothing$; nach unten gilt die duale Bedingung.
Die getrennte Beschränktheit von $S$ und $T$ genügt hierfür
im allgemeinen nicht.
\end{theorem}

\begin{theorem}[Schranken einer Mengenfamilie]
\label{b13guide:thm:family-bounds}
Sei $(S_i)_{i\in I}$ eine beliebige Familie von Teilmengen von $P$.
Dann gilt
\[
 \Lo\Bigl(\bigcup_{i\in I}S_i\Bigr)=\bigcap_{i\in I}\Lo(S_i),\qquad
 \Hi\Bigl(\bigcup_{i\in I}S_i\Bigr)=\bigcap_{i\in I}\Hi(S_i).
\]
Alle Durchschnitte werden innerhalb von $P$ gebildet;
für $I=\varnothing$ ist ihr Wert $P$.
Für den Durchschnitt der Familie gilt
\[
 \bigcup_{i\in I}\Lo(S_i)\subseteq
       \Lo\Bigl(\bigcap_{i\in I}S_i\Bigr),\qquad
 \bigcup_{i\in I}\Hi(S_i)\subseteq
       \Hi\Bigl(\bigcap_{i\in I}S_i\Bigr).
\]

Insbesondere lassen sich die Schranken einer einzelnen Menge als
Durchschnitte ihrer elementweisen Schranken schreiben:
\[
 \Lo(S)=\bigcap_{x\in S}\mathord\downarrow x,\qquad
 \Hi(S)=\bigcap_{x\in S}\mathord\uparrow x.
\]
Auch hier hat ein leerer Durchschnitt den Wert $P$.

\end{theorem}

\begin{theorem}[Unbeschränktheit als Zeugenbedingung]
\label{b13guide:thm:unbounded}
Für $S\subseteq P$ gilt
\[
\begin{aligned}
 \Hi(S)=\varnothing
   &\ \Longleftrightarrow\ \forall u\in P\;\exists x\in S\;(x\nleq u),\\
 \Lo(S)=\varnothing
   &\ \Longleftrightarrow\ \forall \ell\in P\;\exists x\in S\;(\ell\nleq x).
\end{aligned}
\]
Die Bedingung $x\nleq u$ kann nicht allgemein durch $u<x$
ersetzt werden; $x$ und $u$ dürfen unvergleichbar sein.
Dabei bedeutet $u<x$ die Bedingung $u\leq x$ und $u\ne x$.
\end{theorem}
